% !TEX root = ../master_thesis.tex
\chapter{Benchmark Design and Methodology}
\label{chap:methodology}

This chapter defines the benchmark and evaluation protocol used in this thesis.
The central object being scored is the generated UI or Web app artifact. The
leaderboard ranking target is the generator LLM: artifact-level scores are
aggregated across a fixed benchmark item set into model-level leaderboard rows.

The chapter first states the research questions, then defines the benchmark
object model, generation protocol, rendering protocol, metric categories,
dynamic validation strategy, reliability audit, and the mapping from research
questions to methods.

\section{Research Questions}
\label{sec:rqs}

The thesis investigates how to evaluate LLM-generated web interfaces with a
reproducible benchmark and leaderboard. It uses automated scoring, human review,
and LLM-based judging as measurement instruments, but these instruments are not
the main leaderboard target.

\paragraph{RQ1 -- Benchmark and leaderboard design.}
\begin{quote}
\textit{How can a reproducible benchmark and leaderboard be designed for
evaluating LLM-generated web interfaces?}
\end{quote}

RQ1 defines the fixed benchmark dataset, benchmark item schema, generation
protocol, generated artifact format, metadata requirements, and aggregation from
artifact-level scores to generator-model leaderboard rows. Track~A/UIClip-style
material is retained only as a small baseline or sanity check, not as the main
research question or primary contribution.

\paragraph{RQ2 -- Metric operationalization and aggregation.}
\begin{quote}
\textit{How can static technical, static visual, dynamic task-based, and
efficiency metrics be operationalized and aggregated for generated web apps?}
\end{quote}

RQ2 defines the category scores used by the leaderboard: Static Technical
Score, Static Visual Score, Dynamic Score, and Efficiency Score. An optional
Overall Score may be reported, but category-specific rankings remain visible
because different users may prioritize visual quality, functional correctness,
or cost efficiency differently.

\paragraph{RQ3 -- Static and dynamic quality relationship.}
\begin{quote}
\textit{How are static quality scores related to dynamic task-success outcomes
in LLM-generated web apps?}
\end{quote}

RQ3 evaluates the same generated artifacts with both static and dynamic
signals. The analysis reports correlations and mismatch cases, such as visually
strong interfaces that fail tasks or visually plain interfaces that complete
required workflows reliably.

\paragraph{Reliability audit.}
\begin{quote}
\textit{How stable and bias-sensitive are the LLM-based visual or rubric scores
used in the evaluation pipeline?}
\end{quote}

The reliability audit supports the scoring protocol. It records repeated
scoring consistency, order-swap sensitivity for pairwise comparisons when used,
rating-scale sensitivity if feasible, judge-model variance, and invalid or
refusal rates. It is not a separate main research question and does not make
judge-model comparison the thesis target.

\section{Overall Benchmark Design}
\label{sec:design}

The benchmark is based on fixed items and controlled generation. A generator
model receives the same benchmark context and standardized prompt schema as
other models in the same comparison batch. The model produces executable web
artifacts, and those artifacts are scored with the same evaluation protocol.

\textbf{Direct evaluation target.} The generated UI/Web app artifact is scored
for one benchmark item. Scores include technical coverage, visual quality,
dynamic task behavior, and efficiency metadata.

\textbf{Leaderboard ranking target.} The generator LLM or generator model
configuration is ranked. Artifact-level scores are aggregated across benchmark
items into model-level rows.

\textbf{Primary mode.} The main mode is New LLM mode: a generator model
API/interface is added to the pipeline, the fixed benchmark set is run, and the
result is a model row in the leaderboard.

\textbf{Secondary mode.} A single HTML file or zipped project may be scored as
an artifact-level report, but this is secondary. It is not the main leaderboard
basis unless the artifact was generated under the fixed benchmark protocol and
includes complete model, prompt, item, timestamp, and generation metadata.

\begin{figure}[t]
  \centering
  \placeholderfigure{%
    [Figure 3.1: Benchmark pipeline.\\
    Fixed benchmark items and prompt schema feed generator models.\\
    Generated artifacts are scored with static technical, static visual,
    dynamic, and efficiency metrics.\\
    Artifact-level scores aggregate into generator-model leaderboard rows.]}
  \caption{Benchmark and leaderboard pipeline. Generated artifacts are the
  direct scoring target; generator models are the leaderboard ranking target.}
  \label{fig:pipeline}
\end{figure}

\section{Benchmark Item and Generated Artifact}
\label{sec:item-artifact}

A benchmark item is the fixed test specification. It is not a generated output.
Each item contains the information needed to generate and evaluate a web
interface:

\begin{itemize}
  \item natural-language requirement text;
  \item expected pages or routes;
  \item required elements, text, controls, forms, and selectors;
  \item dynamic tasks or workflows;
  \item validation rules;
  \item predefined pages or routes for static visual screenshots;
  \item metadata such as source, category, difficulty, and whether dynamic
  validation is required.
\end{itemize}

The recommended planning schema is:

\begin{verbatim}
benchmark_item/
  requirement.md
  pages.json
  elements.json
  tasks.json
  validation_rules.json
  visual_pages.json
  metadata.json
\end{verbatim}

A generated artifact is the output produced by one generator model for one
benchmark item. The minimal artifact is an executable web implementation plus
generation metadata. The metadata records model identity, provider, prompt ID,
timestamp, provider parameters, token usage, latency, cost, and finish reason.

\section{Dataset Tracks}
\label{sec:tracks}

\subsection{Track A: Baseline or Sanity Check}
\label{sec:track-a}

Track~A uses a reduced UIClip-style screenshot baseline. Its purpose is to
anchor the study in prior screenshot-based GUI evaluation and to provide a
sanity check for pairwise judging. It is not the primary contribution and does
not define the main leaderboard objective.

Track~A can reuse existing UIClip preference labels for a small sample of
screenshots. No new large annotation campaign is required for this baseline.

\subsection{Track B: Generated Web Interfaces}
\label{sec:track-b}

Track~B is the main benchmark track. It uses requirement-driven web interface
generation tasks, primarily from a reduced Vision2Web Level~1/Level~2 subset.
Level~3 full-stack tasks are excluded because they introduce backend state,
deployment, and long-horizon engineering requirements that are outside the
current benchmark scope.

The development prototype should begin with 2--5 benchmark web app items. If
feasible, the thesis experiment may expand to 10--20 benchmark tasks. Expansion
depends on generation cost, scoring cost, rendering reliability, and the amount
of usable dynamic validation available in the selected items.

\section{Generation Protocol}
\label{sec:generation-protocol}

For direct generator comparisons, all models in a comparison batch receive the
same benchmark item context and the same standardized prompt schema. Model-
specific prompt tuning is not allowed within a batch because it would confound
generator quality with prompt engineering.

The main input to the leaderboard pipeline is a generator model API/interface.
For each model, the pipeline loads the fixed benchmark items, applies the
standardized prompt, saves generated artifacts, records generation metadata, and
runs the evaluation protocol.

The evaluation should avoid presenting any fixed truncating output-token cap as
the final policy. If a provider requires a maximum output setting, it should be
set high enough to allow complete outputs where feasible. Actual input tokens,
output tokens, total tokens, finish reason, cost, and latency are recorded as
efficiency and reproducibility data.

Prompts that affect thesis results, metric computation, benchmark protocol,
human annotation, output schemas, or reproducibility are versioned and preserved
according to the prompt preservation policy.

\section{Rendering and Screenshot Protocol}
\label{sec:rendering}

Generated artifacts are rendered locally with a controlled browser environment.
Rendering produces the evidence used by later scoring modules.

For static visual scoring, the pipeline visits only the predefined pages or
routes specified by the benchmark item and captures standardized screenshots.
The browser engine, viewport, zoom level, wait rule, and screenshot mode are
kept fixed across artifacts. These screenshots are the input to static visual
evaluation.

Agent trace screenshots are not used as the main static visual scoring input.
They may be stored for debugging and dynamic failure analysis only, because
agent traces can vary in length, path, and content across generated apps.

Rendering failures such as blank pages, fatal JavaScript errors that prevent
visible content, or missing entry files are recorded as artifact-level failures.

\section{Metric Categories}
\label{sec:metric-categories}

\subsection{Static Technical Score}

Static technical metrics evaluate structure and requirement coverage without
executing a user task. Candidate submetrics include HTML completeness, required
element coverage, required text coverage, route declaration coverage, link
target declaration, form field coverage, selector or ID coverage, accessibility
basics, and static requirement fidelity.

Static technical checks can verify that a required button exists and declares a
target. They do not prove that clicking the button completes a task; that is a
dynamic metric.

\subsection{Static Visual Score}

Static visual metrics evaluate standardized predefined-route screenshots.
Candidate visual dimensions include visual hierarchy, information organization,
typography and readability, contrast, spacing and alignment, consistency, and
overall aesthetic quality.

Page-level visual scores are aggregated into an artifact-level Static Visual
Score. Agent trace screenshots are excluded from the main static visual score.

\subsection{Dynamic Score}

Dynamic metrics evaluate interaction behavior. A metric is dynamic when it
requires clicking, typing, navigating, submitting, waiting for state changes, or
validating task completion.

Candidate submetrics include task success rate, step success rate, route
execution success, button functionality success, form completion success,
state-change success, retry count, failure localization, and robustness or
pass@k.

The current implementation may use lightweight workflow or action-plan
validation. A future implementation may replace the executor with a real
browser-based computer-use agent, but the metric schema should remain stable:
change the executor, not the metric schema.

\subsection{Efficiency Score}

Efficiency metrics track input tokens, output tokens, total tokens, generation
cost, evaluation cost, generation latency, evaluation latency, and
cost-performance trade-offs.

Efficiency is reported both as raw metadata and, when a reference is defined, as
a normalized score. Cost-performance plots and Pareto frontiers are used to show
whether a model is dominated by another model with both lower cost and higher
quality.

\subsection{Optional Overall Score}

The initial optional composite score is:

\begin{verbatim}
Overall Score =
0.25 * Static Technical Score
+ 0.25 * Static Visual Score
+ 0.35 * Dynamic Score
+ 0.15 * Efficiency Score
\end{verbatim}

This composite is a reporting convenience. Category-specific scores remain
visible and are interpreted separately.

\section{Dynamic Validation Strategy}
\label{sec:dynamic-method}

The thesis first implements lightweight dynamic validation rather than heavy
real-agent execution. This choice keeps the pipeline reproducible and avoids
confounding generated-interface quality with the planning or grounding failures
of a separate agent.

For Vision2Web-derived Level~2 items, dynamic tasks can be derived from
workflow actions and validations. The evaluator checks whether the generated
artifact supports the expected route, form, state, or content outcome. Current
route or workflow simulation results are treated as lightweight dynamic
validation, not as a full real-agent usability study.

Real-agent execution remains an optional future upgrade or small cross-check
after the main pipeline is stable. If added later, the same task success and
failure taxonomy should be reused so that only the executor changes.

\section{Failure Taxonomy}
\label{sec:failure-taxonomy}

Failures are recorded so that leaderboard scores remain interpretable:

\begin{itemize}
  \item \textbf{Generation failure}: no artifact, provider timeout, invalid
  response, or truncated unusable output.
  \item \textbf{Render failure}: artifact exists but cannot be rendered into
  visible content.
  \item \textbf{Static technical failure}: required structure, route, element,
  text, selector, or basic accessibility evidence is missing.
  \item \textbf{Static visual failure}: standardized screenshot shows severe
  visual defects such as overlap, unreadable text, or unusable layout.
  \item \textbf{Dynamic failure}: required task, route, form, button, or state
  change fails under the dynamic validation protocol.
  \item \textbf{Evaluator/tool failure}: scoring or validation tool fails for
  reasons unrelated to the generated artifact.
\end{itemize}

Evaluator or tool failures are reported separately from generated-artifact
quality when possible.

\section{Reliability Audit}
\label{sec:reliability-audit}

LLM-based visual or rubric scoring can be affected by prompt sensitivity,
position bias, scale effects, model variance, and invalid outputs. These issues
are audited but are not the main thesis target.

The audit may include repeated scoring consistency, order-swapped pairwise
comparisons where pairwise scoring is used, rating-scale sensitivity if
feasible, judge model variance, and invalid or refusal rate. Audit results are
reported as scoring reliability evidence and threats to validity.

\section{Research Question -- Method Map}
\label{sec:rq-map}

Table~\ref{tab:rq-map} maps the research questions to data, methods, metrics,
and outputs.

\begin{table}[t]
  \centering
  \caption{Research question -- method mapping.}
  \label{tab:rq-map}
  \small
  \setlength{\tabcolsep}{4pt}
  \begin{tabular}{p{0.7cm} p{4.0cm} p{2.4cm} p{3.0cm} p{3.9cm}}
    \toprule
    \textbf{RQ} & \textbf{Question} & \textbf{Data} & \textbf{Method}
      & \textbf{Primary outputs} \\
    \midrule
    1 & Benchmark and leaderboard design
      & Fixed benchmark items; generated artifacts
      & New LLM mode; artifact scoring; model-level aggregation
      & Benchmark schema; prompt protocol; leaderboard row format \\
    \addlinespace
    2 & Metric operationalization and aggregation
      & Generated artifacts; screenshots; task validations; metadata
      & Static technical, static visual, dynamic, and efficiency scoring
      & Category scores; optional Overall Score; cost-performance plots \\
    \addlinespace
    3 & Static and dynamic quality relationship
      & Same Track~B artifacts scored statically and dynamically
      & Correlation analysis and mismatch case analysis
      & Static-dynamic correlations; qualitative divergence cases \\
    \addlinespace
    Audit & Scoring reliability
      & Selected scoring outputs
      & Repeated scoring, order-swap checks, variance and invalid-rate logs
      & Reliability and bias audit results \\
    \bottomrule
  \end{tabular}
\end{table}
